\section{Field Surface and Volume Inference (FiSure, ``Eitri'')}

    \subsection{Motivation and Concept}

        To compare subfields or regions within Tensor Omics space, we often need 
        to quantify their \textit{geometric extent} --- the shape, surface, and volume 
        occupied by a field of vectors.  
        The \textbf{Field Surface and Volume Inference (FiSure)} method provides a 
        model-free procedure to estimate these quantities directly from the coordinates 
        and projections of field vectors.  

        Conceptually, FiSure treats a field as a set of vectors 
        \( \{\vec{x}_i\} \) distributed around a \textit{semantic axis} or 
        \textit{handle} (e.g. a Mjölnir manifold), and reconstructs the 
        outer surface that encloses them.  
        The method is analogous to measuring the radial extent of a 
        magnetic flux tube or a fluid jet along its main trajectory.

    \subsection{Algorithmic Outline}

        Assume a learned one-dimensional semantic axis 
        \(\vec{h}(t) \in \mathbb{R}^n\) (the \textbf{handle}) representing 
        the main trajectory of the field.  
        For a field consisting of coordinates 
        \( \{\vec{x}_i\}_{i=1}^{m}\), FiSure proceeds as follows:

        \begin{enumerate}
            \item \textbf{Projection onto the Handle.}  
                  Each field point is projected onto the handle by nearest-point projection:
                  \[
                      t_i = \arg\min_t \| \vec{x}_i - \vec{h}(t) \|.
                  \]
                  This yields the axial coordinate \(t_i\) of each field vector.

            \item \textbf{Handle Segmentation.}  
                  The handle is divided into \(k\) equal-percentile segments 
                  \( \{[t_j, t_{j+1}]\} \), such that each segment corresponds to 
                  an equal fraction of projected vectors.  
                  Each segment is approximated by a straight line between its 
                  endpoints \(\vec{h}(t_j)\) and \(\vec{h}(t_{j+1})\).

            \item \textbf{Local Coordinate Frame.}  
                  For each segment, a local cylindrical coordinate system is constructed:
                  \[
                      \text{axis vector: } \vec{a}_j = \frac{\vec{h}(t_{j+1}) - \vec{h}(t_j)}
                                          {\|\vec{h}(t_{j+1}) - \vec{h}(t_j)\|}.
                  \]
                  Orthogonal basis vectors are defined as 
                  \( \{\vec{u}_j, \vec{v}_j\} \), spanning the plane perpendicular to \(\vec{a}_j\).

            \item \textbf{Angular Sampling.}  
                  The circular cross-section around the segment is divided into 
                  \(c\) angular bins:
                  \[
                      \theta_k = \frac{2\pi k}{c}, \quad k = 1,\ldots,c.
                  \]
                  A fixed global reference axis (e.g. the first axis of Tensor Omics space) 
                  is used to align the angular bins between segments.

            \item \textbf{Radial Extent and Surface Extraction.}  
                  For each angular bin, FiSure computes the radial distance of all 
                  vectors whose projections fall within the segment and angle range:
                  \[
                      r_{j,k} = \max_{\vec{x}_i \in \mathcal{S}_{j,k}} 
                      \| (\vec{x}_i - \vec{h}(t_j)) - 
                      \langle \vec{x}_i - \vec{h}(t_j), \vec{a}_j \rangle \vec{a}_j \|.
                  \]
                  The point attaining this maximum is recorded as the local surface point:
                  \[
                      \vec{p}_{j,k} = \vec{h}(t_j) + r_{j,k} \,
                           (\cos\theta_k \, \vec{u}_j + \sin\theta_k \, \vec{v}_j).
                  \]
                  The corresponding outward normal vector is
                  \[
                      \vec{n}_{j,k} = 
                      \frac{\vec{p}_{j,k} - \vec{h}(t_j)}
                      {\|\vec{p}_{j,k} - \vec{h}(t_j)\|}.
                  \]

            \item \textbf{Surface and Volume Assembly.}  
                The set of surface points 
                \(\mathcal{P} = \{\vec{p}_{j,k}\} \) defines a discrete 
                surface mesh.  
                The local volume of the segment between 
                \(t_j\) and \(t_{j+1}\) can be approximated as a truncated cylinder:
                \[
                    V_j \approx 
                    \frac{\pi \, \Delta s_j}{c}
                    \sum_{k=1}^{c}
                    \frac{r_{j,k}^{2} + r_{j+1,k}^{2}}{2},
                \]
                where \(\Delta s_j = \|\vec{h}(t_{j+1}) - \vec{h}(t_j)\|\).
                Summing over all segments yields the total field volume:
                \[
                    V_{\text{field}} = \sum_{j=1}^{k-1} V_j.
                \]
        \end{enumerate}

    \subsection{Interpretation of FiSure Quantities}

        FiSure yields several interpretable geometric measures:

        \begin{itemize}
            \item \textbf{Surface Vectors:}  
                  Each \((\vec{p}_{j,k}, \vec{n}_{j,k})\) pair represents a point on the 
                  field’s boundary and its outward orientation.  
                  These can be analyzed as a secondary \textit{surface vector field}, 
                  quantifying how the field opens, closes, or bends in Tensor space.

            \item \textbf{Local Volume Elements:}  
                  The segment volumes \(V_j\) allow normalization of other descriptors, 
                  such as energy or rank, by the spatial size of the field.

            \item \textbf{Field Morphology:}  
                  Variations in \(r_{j,k}\) and curvature of the handle describe 
                  anisotropy and torsion of the field.  
                  A stable \(r_{j,k}\) profile implies a cylindrical, consistent field; 
                  strong local expansion or contraction indicates field divergence or collapse.
        \end{itemize}

    \subsection{Applications and Interpretation}

        FiSure provides a direct geometric connection between 
        field structure and function.  
        In the context of Tensor Omics, examples include:

        \begin{itemize}
            \item Measuring the \textbf{semantic volume} occupied by gene families 
                  associated with aging or stress responses.
            \item Quantifying whether such semantic volumes 
                  \textbf{expand, contract, or twist} along a 
                  developmental or disease axis.
            \item Identifying rotational or asymmetrical \textbf{surface deformations} 
                  that indicate functional bifurcation (e.g., 
                  tissue-specific reprogramming).
        \end{itemize}

        The same principles apply to physical or environmental fields:
        for a wind field, FiSure reconstructs the shape of a coherent 
        airflow tube; for a magnetic field, it delineates regions of 
        strong magnetic flux concentration.

    \subsection{Summary}

        FiSure extends Tensor Omics by providing a quantitative 
        method for reconstructing the \textbf{shape} and 
        \textbf{extent} of a field.  
        By defining surfaces and volumes empirically from vector distributions, 
        it bridges local field descriptors (energy, rank, rotation) with 
        global geometric form, completing the description of 
        high-dimensional semantic dynamics in Tensor space.
